\thispagestyle{empty}
\begin{center}
    \Huge{\textbf{我掌握的流体力学}}\\
    \vspace{1.5cm}
    \Huge\textbf{前言}
\end{center}

这份笔记整理始于2026年端午节，初衷是为推免申请做准备——希望在有限时间内，系统梳理流体力学
相关的核心概念，便于向老师系统展示我在本科前三年掌握的流体力学相关概念。在整理过程中，
复变函数、矢量恒等式等数学工具默认已在先修课程中掌握，故仅简要提及推导思路或者直接使用。
预计九月完成全部整理，涵盖流体力学基础（包括实验流体力学）、气体动力学、连续介质力学、理论力学，
以及热力学与统计物理等方向。目前正在撰写第一部分——流体力学基础。
\href{https://github.com/nobodyMa/nobodyMa_note_review_Fluiddynamics/blob/main/%E6%88%91%E6%8E%8C%E6%8F%A1%E7%9A%84%E6%B5%81%E4%BD%93%E5%8A%9B%E5%AD%A6.pdf}{最新版本的笔记}将持续更新在GitHub上，
欢迎大家提出建议。
\vspace{0.5cm}

\begin{flushright}
    作者：无名氏马\\
    \today
\end{flushright}

\newpage                      % 新的一页
\pagestyle{plain}             % 设置页眉和页脚的排版方式（plain：页眉是空的，页脚只包含一个居中的页码）
\setcounter{page}{1}          % 重新定义页码从第一页开始
\pagenumbering{Roman}         % 使用大写的罗马数字作为页码
\tableofcontents              % 生成目录

\newpage                      % 以下是正文
\pagestyle{plain}
\setcounter{page}{1}          % 使用阿拉伯数字作为页码
\pagenumbering{arabic}
% \setcounter{chapter}{-1}    % 设置 -1 可作为第零章绪论从第零章开始
